
\documentclass{article}
\usepackage{colortbl}
\usepackage{makecell}
\usepackage{multirow}
\usepackage{supertabular}

\begin{document}

\newcounter{utterance}

\centering \large Interaction Transcript for game `clean\_up', experiment `1\_medium\_5obj\_sk', episode 1 with gemma{-}4{-}26b{-}a4b{-}it.
\vspace{24pt}

{ \footnotesize  \setcounter{utterance}{1}
\setlength{\tabcolsep}{0pt}
\begin{supertabular}{c@{$\;$}|p{.15\linewidth}@{}p{.15\linewidth}p{.15\linewidth}p{.15\linewidth}p{.15\linewidth}p{.15\linewidth}}
    \# & $\;$A & \multicolumn{4}{c}{Game Master} & $\;\:$B\\
    \hline

    \theutterance \stepcounter{utterance}  
    & & \multicolumn{4}{p{0.6\linewidth}}{
        \cellcolor[rgb]{0.9,0.9,0.9}{
            \makecell[{{p{\linewidth}}}]{
                \texttt{\tiny{[P1$\langle$GM]}}
                \texttt{Som váš vedúci hry a hráte spoločnú hru s nasledujúcou mriežkou ako hracím plánom:} \\
\\ 
\texttt{```} \\
\texttt{ 1234567} \\
\texttt{╔══╤══╤═╗} \\
\texttt{║◌◌│T◌│◌║ 1} \\
\texttt{╟──┼──┴─╢ 2} \\
\texttt{║WI│◌◌◌◌║ 3} \\
\texttt{╟──┘◌◌◌◌║ 4} \\
\texttt{║◌◌◌H◌◌C║ 5} \\
\texttt{║◌◌┌─┬──╢ 6} \\
\texttt{║◌◌│◌│◌◌║ 7} \\
\texttt{╚══╧═╧══╝} \\
\\ 
\texttt{```} \\
\\ 
\texttt{* Horný okraj zobrazuje súradnice x rastúce doprava a pravý okraj súradnice y rastúce nadol.} \\
\texttt{* Na vašej mriežke sú náhodne umiestnené nasledujúce objekty: 'W', 'I', 'T', 'C', 'H'.} \\
\\ 
\\ 
\texttt{Druhý hráč vidí variant hracej plochy, kde sú objekty umiestnené na rôznych náhodných miestach. Vy nevidíte hraciu plochu druhého hráča a on nevidí tú vašu.} \\
\\ 
\\ 
\texttt{**Cieľ:**} \\
\\ 
\texttt{Obaja hráči musia presunúť objekty na svojom pozadí tak, aby sa rovnaké objekty ocitli na rovnakých súradniciach. Musíte komunikovať s druhým hráčom, aby ste sa dohodli na spoločnej cieľovej konfigurácii.} \\
\\ 
\\ 
\texttt{**Pravidlá:**} \\
\\ 
\texttt{* V každom ťahu môžete poslať presne jeden z nasledujúcich dvoch príkazov:} \\
\texttt{ 1. `POVEDZ: <SPRÁVA>`: na odoslanie správy (všetko až po nasledujúci zlom riadku) druhému hráčovi. Prepošlem ju vášmu partnerovi.} \\
\texttt{2. `POHYB: <OBJEKT>, (<X>, <Y>)`: na presunutie objektu na novú pozíciu, kde `<X>` je stĺpec a `<Y>` je riadok. Budem vás informovať, či bol váš presun úspešný alebo nie.} \\
\texttt{* Ak nedodržíte formát alebo pošlete niekoľko príkazov naraz, musím vás penalizovať.} \\
\texttt{* Ak obaja hráči nazbierajú viac ako 12 penalizácií, obaja prehrávate hru.} \\
\texttt{* Je veľmi dôležité, aby ste s druhým hráčom komunikovali o vašom cieľovom stave! *Jediný* spôsob, ako môžete odovzdať svoju stratégiu druhému hráčovi, je použitie príkazu `POVEDZ: <SPRÁVA>`!} \\
\\ 
\\ 
\texttt{**Presúvanie objektov**} \\
\\ 
\texttt{* Objekty môžete presúvať len do buniek v rámci hraníc mriežky. Cieľová bunka musí byť prázdna, t. j. musí obsahovať iba symbol '◌'.} \\
\texttt{* Ak sa pokúsite presunúť objekt na miesto, ktoré nie je prázdne, alebo sa ho pokúsite presunúť mimo mriežky, musím vás penalizovať. Dostanete ďalší pokus.} \\
\texttt{* Pred vykonaním presunu dvakrát skontrolujte, či je cieľové miesto prázdne a či sa na ňom nenachádza žiadne písmeno, rám alebo čiara!} \\
\\ 
\\ 
\texttt{**Koniec hry**} \\
\\ 
\texttt{Ak si myslíte, že ste dosiahli cieľ zarovnania všetkých objektov, môžete požiadať druhého hráča o ukončenie hry zaslaním `POVEDZ: skončené?`. Ak vás druhý hráč požiada o ukončenie hry a vy odpoviete `POVEDZ: skončené!`, hra sa skončí.} \\
\\ 
\texttt{Obaja hráči vyhrajú, ak hra skončí do 20 kôl, pričom jedno kolo je definované ako odoslanie platného príkazu každým z dvoch hráčov.} \\
\\ 
\texttt{**Bodovanie:**} \\
\\ 
\texttt{Čím bližšie sú rovnaké objekty na oboch hracích plochách, tým viac bodov získate. Penalizácie znižujú počet vašich bodov. Podarí sa vám prekonať rekord?} \\
\\ 
\\ 
\texttt{Pošlite správu druhému hráčovi, aby ste začali hru!} \\
            }
        }
    }
    & \\ \\

    \theutterance \stepcounter{utterance}  
    & \multicolumn{4}{p{0.6\linewidth}}{
        \cellcolor[rgb]{0.9,0.9,0.9}{
            \makecell[{{p{\linewidth}}}]{
                \texttt{\tiny{[P1$\rangle$GM]}}
                \texttt{POVEDZ: Ahoj! Moja súčasná konfigurácia je: W(1,3), I(2,3), T(4,1), H(4,5), C(6,5). Navrhujem, aby sme sa pokúsili zarovnať všetky objekty do stredovej časti mriežky alebo do konkrétnych súradníc, ktoré sú pre oboch nás ľahko dostupné. Aké máš ty objekty a akú navrhuješ cieľovú konfiguráciu?} \\
            }
        }
    }
    & & \\ \\

    \theutterance \stepcounter{utterance}  
    & & & \multicolumn{2}{p{0.3\linewidth}}{
        \cellcolor[rgb]{0.9,0.9,0.9}{
            \makecell[{{p{\linewidth}}}]{
                \texttt{\tiny{[GM$|$GM]}}
                \texttt{POVEDZ: Ahoj! Moja súčasná konfigurácia je: W(1,3), I(2,3), T(4,1), H(4,5), C(6,5). Navrhujem, aby sme sa pokúsili zarovnať všetky objekty do stredovej časti mriežky alebo do konkrétnych súradníc, ktoré sú pre oboch nás ľahko dostupné. Aké máš ty objekty a akú navrhuješ cieľovú konfiguráciu?} \\
            }
        }
    }
    & & \\ \\

    \theutterance \stepcounter{utterance}  
    & & \multicolumn{4}{p{0.6\linewidth}}{
        \cellcolor[rgb]{0.9,0.9,0.9}{
            \makecell[{{p{\linewidth}}}]{
                \texttt{\tiny{[GM$\rangle$P2]}}
                \texttt{Som váš vedúci hry a hráte spoločnú hru s nasledujúcou mriežkou ako hracím plánom:} \\
\\ 
\texttt{```} \\
\texttt{ 1234567} \\
\texttt{╔══╤══╤═╗} \\
\texttt{║◌◌│W◌│I║ 1} \\
\texttt{╟──┼──┴─╢ 2} \\
\texttt{║◌◌│◌◌◌◌║ 3} \\
\texttt{╟──┘T◌◌◌║ 4} \\
\texttt{║H◌◌◌◌◌◌║ 5} \\
\texttt{║◌◌┌─┬──╢ 6} \\
\texttt{║C◌│◌│◌◌║ 7} \\
\texttt{╚══╧═╧══╝} \\
\\ 
\texttt{```} \\
\\ 
\texttt{* Horný okraj zobrazuje súradnice x rastúce doprava a pravý okraj súradnice y rastúce nadol.} \\
\texttt{* Na vašej mriežke sú náhodne umiestnené nasledujúce objekty: 'W', 'I', 'T', 'C', 'H'.} \\
\\ 
\\ 
\texttt{Druhý hráč vidí variant hracej plochy, kde sú objekty umiestnené na rôznych náhodných miestach. Vy nevidíte hraciu plochu druhého hráča a on nevidí tú vašu.} \\
\\ 
\\ 
\texttt{**Cieľ:**} \\
\\ 
\texttt{Obaja hráči musia presunúť objekty na svojom pozadí tak, aby sa rovnaké objekty ocitli na rovnakých súradniciach. Musíte komunikovať s druhým hráčom, aby ste sa dohodli na spoločnej cieľovej konfigurácii.} \\
\\ 
\\ 
\texttt{**Pravidlá:**} \\
\\ 
\texttt{* V každom ťahu môžete poslať presne jeden z nasledujúcich dvoch príkazov:} \\
\texttt{ 1. `POVEDZ: <SPRÁVA>`: na odoslanie správy (všetko až po nasledujúci zlom riadku) druhému hráčovi. Prepošlem ju vášmu partnerovi.} \\
\texttt{2. `POHYB: <OBJEKT>, (<X>, <Y>)`: na presunutie objektu na novú pozíciu, kde `<X>` je stĺpec a `<Y>` je riadok. Budem vás informovať, či bol váš presun úspešný alebo nie.} \\
\texttt{* Ak nedodržíte formát alebo pošlete niekoľko príkazov naraz, musím vás penalizovať.} \\
\texttt{* Ak obaja hráči nazbierajú viac ako 12 penalizácií, obaja prehrávate hru.} \\
\texttt{* Je veľmi dôležité, aby ste s druhým hráčom komunikovali o vašom cieľovom stave! *Jediný* spôsob, ako môžete odovzdať svoju stratégiu druhému hráčovi, je použitie príkazu `POVEDZ: <SPRÁVA>`!} \\
\\ 
\\ 
\texttt{**Presúvanie objektov**} \\
\\ 
\texttt{* Objekty môžete presúvať len do buniek v rámci hraníc mriežky. Cieľová bunka musí byť prázdna, t. j. musí obsahovať iba symbol '◌'.} \\
\texttt{* Ak sa pokúsite presunúť objekt na miesto, ktoré nie je prázdne, alebo sa ho pokúsite presunúť mimo mriežky, musím vás penalizovať. Dostanete ďalší pokus.} \\
\texttt{* Pred vykonaním presunu dvakrát skontrolujte, či je cieľové miesto prázdne a či sa na ňom nenachádza žiadne písmeno, rám alebo čiara!} \\
\\ 
\\ 
\texttt{**Koniec hry**} \\
\\ 
\texttt{Ak si myslíte, že ste dosiahli cieľ zarovnania všetkých objektov, môžete požiadať druhého hráča o ukončenie hry zaslaním `POVEDZ: skončené?`. Ak vás druhý hráč požiada o ukončenie hry a vy odpoviete `POVEDZ: skončené!`, hra sa skončí.} \\
\\ 
\texttt{Obaja hráči vyhrajú, ak hra skončí do 20 kôl, pričom jedno kolo je definované ako odoslanie platného príkazu každým z dvoch hráčov.} \\
\\ 
\texttt{**Bodovanie:**} \\
\\ 
\texttt{Čím bližšie sú rovnaké objekty na oboch hracích plochách, tým viac bodov získate. Penalizácie znižujú počet vašich bodov. Podarí sa vám prekonať rekord?} \\
\\ 
\texttt{Druhý hráč začal hru odoslaním tejto správy:} \\
\\ 
\texttt{"Ahoj! Moja súčasná konfigurácia je: W(1,3), I(2,3), T(4,1), H(4,5), C(6,5). Navrhujem, aby sme sa pokúsili zarovnať všetky objekty do stredovej časti mriežky alebo do konkrétnych súradníc, ktoré sú pre oboch nás ľahko dostupné. Aké máš ty objekty a akú navrhuješ cieľovú konfiguráciu?"} \\
\\ 
\texttt{Aký je váš prvý príkaz?} \\
            }
        }
    }
    & \\ \\

    \theutterance \stepcounter{utterance}  
    & & & \multicolumn{4}{p{0.6\linewidth}}{
        \cellcolor[rgb]{0.9,0.9,0.9}{
            \makecell[{{p{\linewidth}}}]{
                \texttt{\tiny{[GM$\langle$P2]}}
                \texttt{POVEDZ: Ahoj! Moja konfigurácia je: W(4,1), I(7,1), T(4,4), H(1,5), C(1,7). Navrhujem, aby sme sa zamerali na jednoduchú konfiguráciu v strednej časti. Skúsme zarovnať objekty na tieto súradnice: W(4,2), I(5,2), T(4,3), H(3,3), C(2,3). Čo na to povieš?} \\
            }
        }
    }
     \\ \\

    \theutterance \stepcounter{utterance}  
    & & & \multicolumn{2}{p{0.3\linewidth}}{
        \cellcolor[rgb]{0.9,0.9,0.9}{
            \makecell[{{p{\linewidth}}}]{
                \texttt{\tiny{[GM$|$GM]}}
                \texttt{POVEDZ: Ahoj! Moja konfigurácia je: W(4,1), I(7,1), T(4,4), H(1,5), C(1,7). Navrhujem, aby sme sa zamerali na jednoduchú konfiguráciu v strednej časti. Skúsme zarovnať objekty na tieto súradnice: W(4,2), I(5,2), T(4,3), H(3,3), C(2,3). Čo na to povieš?} \\
            }
        }
    }
    & & \\ \\

    \theutterance \stepcounter{utterance}  
    & & \multicolumn{4}{p{0.6\linewidth}}{
        \cellcolor[rgb]{0.9,0.9,0.9}{
            \makecell[{{p{\linewidth}}}]{
                \texttt{\tiny{[P1$\langle$GM]}}
                \texttt{Vaša správa bola odovzdaná druhému hráčovi.} \\
\\ 
\\ 
\texttt{Práve hráte kolo 1 z maximálne 20.} \\
\texttt{Spoločne ste nazbierali 0 z 12 pokút.} \\
\texttt{Druhý hráč poslal túto správu:} \\
\\ 
\\ 
\texttt{"Ahoj! Moja konfigurácia je: W(4,1), I(7,1), T(4,4), H(1,5), C(1,7). Navrhujem, aby sme sa zamerali na jednoduchú konfiguráciu v strednej časti. Skúsme zarovnať objekty na tieto súradnice: W(4,2), I(5,2), T(4,3), H(3,3), C(2,3). Čo na to povieš?"} \\
\\ 
\\ 
\texttt{Aký je váš ďalší príkaz?} \\
            }
        }
    }
    & \\ \\

    \theutterance \stepcounter{utterance}  
    & \multicolumn{4}{p{0.6\linewidth}}{
        \cellcolor[rgb]{0.9,0.9,0.9}{
            \makecell[{{p{\linewidth}}}]{
                \texttt{\tiny{[P1$\rangle$GM]}}
                \texttt{POVEDZ: To je dobrý plán! Tieto súradnice sú v strednej časti a vyzerajú na to, že sú pre nás oboch dostupné. Prijímam tvoj návrh: W(4,2), I(5,2), T(4,3), H(3,3), C(2,3). Začnem s presunom mojich objektov. Najprv presunú T na novú pozíciu.} \\
            }
        }
    }
    & & \\ \\

    \theutterance \stepcounter{utterance}  
    & & & \multicolumn{2}{p{0.3\linewidth}}{
        \cellcolor[rgb]{0.9,0.9,0.9}{
            \makecell[{{p{\linewidth}}}]{
                \texttt{\tiny{[GM$|$GM]}}
                \texttt{POVEDZ: To je dobrý plán! Tieto súradnice sú v strednej časti a vyzerajú na to, že sú pre nás oboch dostupné. Prijímam tvoj návrh: W(4,2), I(5,2), T(4,3), H(3,3), C(2,3). Začnem s presunom mojich objektov. Najprv presunú T na novú pozíciu.} \\
            }
        }
    }
    & & \\ \\

    \theutterance \stepcounter{utterance}  
    & & \multicolumn{4}{p{0.6\linewidth}}{
        \cellcolor[rgb]{0.9,0.9,0.9}{
            \makecell[{{p{\linewidth}}}]{
                \texttt{\tiny{[GM$\rangle$P2]}}
                \texttt{Vaša správa bola odovzdaná druhému hráčovi.} \\
\\ 
\\ 
\texttt{Práve hráte kolo 2 z maximálne 20.} \\
\texttt{Spoločne ste nazbierali 0 z 12 pokút.} \\
\texttt{Druhý hráč poslal túto správu:} \\
\\ 
\\ 
\texttt{"To je dobrý plán! Tieto súradnice sú v strednej časti a vyzerajú na to, že sú pre nás oboch dostupné. Prijímam tvoj návrh: W(4,2), I(5,2), T(4,3), H(3,3), C(2,3). Začnem s presunom mojich objektov. Najprv presunú T na novú pozíciu."} \\
\\ 
\\ 
\texttt{Aký je váš ďalší príkaz?} \\
            }
        }
    }
    & \\ \\

    \theutterance \stepcounter{utterance}  
    & & & \multicolumn{4}{p{0.6\linewidth}}{
        \cellcolor[rgb]{0.9,0.9,0.9}{
            \makecell[{{p{\linewidth}}}]{
                \texttt{\tiny{[GM$\langle$P2]}}
                \texttt{POHYB: T, (4, 3)} \\
            }
        }
    }
     \\ \\

    \theutterance \stepcounter{utterance}  
    & & & \multicolumn{2}{p{0.3\linewidth}}{
        \cellcolor[rgb]{0.9,0.9,0.9}{
            \makecell[{{p{\linewidth}}}]{
                \texttt{\tiny{[GM$|$GM]}}
                \texttt{POHYB: T, (4, 3)} \\
            }
        }
    }
    & & \\ \\

    \theutterance \stepcounter{utterance}  
    & & & \multicolumn{2}{p{0.3\linewidth}}{
        \cellcolor[rgb]{0.9,0.9,0.9}{
            \makecell[{{p{\linewidth}}}]{
                \texttt{\tiny{[GM$|$GM]}}
                \texttt{T moved to (4, 3)} \\
            }
        }
    }
    & & \\ \\

    \theutterance \stepcounter{utterance}  
    & & \multicolumn{4}{p{0.6\linewidth}}{
        \cellcolor[rgb]{0.9,0.9,0.9}{
            \makecell[{{p{\linewidth}}}]{
                \texttt{\tiny{[P1$\langle$GM]}}
                \texttt{Vaša správa bola odovzdaná druhému hráčovi.} \\
\\ 
\\ 
\texttt{Práve hráte kolo 2 z maximálne 20.} \\
\texttt{Spoločne ste nazbierali 0 z 12 pokút.} \\
\texttt{Druhý hráč presunul objekt na svojej mriežke.} \\
\\ 
\\ 
\texttt{Aký je váš ďalší príkaz?} \\
            }
        }
    }
    & \\ \\

    \theutterance \stepcounter{utterance}  
    & \multicolumn{4}{p{0.6\linewidth}}{
        \cellcolor[rgb]{0.9,0.9,0.9}{
            \makecell[{{p{\linewidth}}}]{
                \texttt{\tiny{[P1$\rangle$GM]}}
                \texttt{POHYB: T, (4, 3)} \\
            }
        }
    }
    & & \\ \\

    \theutterance \stepcounter{utterance}  
    & & & \multicolumn{2}{p{0.3\linewidth}}{
        \cellcolor[rgb]{0.9,0.9,0.9}{
            \makecell[{{p{\linewidth}}}]{
                \texttt{\tiny{[GM$|$GM]}}
                \texttt{POHYB: T, (4, 3)} \\
            }
        }
    }
    & & \\ \\

    \theutterance \stepcounter{utterance}  
    & & & \multicolumn{2}{p{0.3\linewidth}}{
        \cellcolor[rgb]{0.9,0.9,0.9}{
            \makecell[{{p{\linewidth}}}]{
                \texttt{\tiny{[GM$|$GM]}}
                \texttt{T moved to (4, 3)} \\
            }
        }
    }
    & & \\ \\

    \theutterance \stepcounter{utterance}  
    & & \multicolumn{4}{p{0.6\linewidth}}{
        \cellcolor[rgb]{0.9,0.9,0.9}{
            \makecell[{{p{\linewidth}}}]{
                \texttt{\tiny{[GM$\rangle$P2]}}
                \texttt{Úspešne presunutý 'T' na (4,3). Vaša aktualizovaná mriežka vyzerá takto:} \\
\texttt{```} \\
\texttt{ 1234567} \\
\texttt{╔══╤══╤═╗} \\
\texttt{║◌◌│W◌│I║ 1} \\
\texttt{╟──┼──┴─╢ 2} \\
\texttt{║◌◌│T◌◌◌║ 3} \\
\texttt{╟──┘◌◌◌◌║ 4} \\
\texttt{║H◌◌◌◌◌◌║ 5} \\
\texttt{║◌◌┌─┬──╢ 6} \\
\texttt{║C◌│◌│◌◌║ 7} \\
\texttt{╚══╧═╧══╝} \\
\texttt{```} \\
\texttt{Práve hráte kolo 3 z maximálne 20.} \\
\texttt{Spoločne ste nazbierali 0 z 12 pokút.} \\
\texttt{Druhý hráč presunul objekt na svojej mriežke.} \\
\\ 
\\ 
\texttt{Aký je váš ďalší príkaz?} \\
            }
        }
    }
    & \\ \\

    \theutterance \stepcounter{utterance}  
    & & & \multicolumn{4}{p{0.6\linewidth}}{
        \cellcolor[rgb]{0.9,0.9,0.9}{
            \makecell[{{p{\linewidth}}}]{
                \texttt{\tiny{[GM$\langle$P2]}}
                \texttt{POHYB: W, (4, 2)} \\
            }
        }
    }
     \\ \\

    \theutterance \stepcounter{utterance}  
    & & & \multicolumn{2}{p{0.3\linewidth}}{
        \cellcolor[rgb]{0.9,0.9,0.9}{
            \makecell[{{p{\linewidth}}}]{
                \texttt{\tiny{[GM$|$GM]}}
                \texttt{POHYB: W, (4, 2)} \\
            }
        }
    }
    & & \\ \\

    \theutterance \stepcounter{utterance}  
    & & & \multicolumn{2}{p{0.3\linewidth}}{
        \cellcolor[rgb]{0.9,0.9,0.9}{
            \makecell[{{p{\linewidth}}}]{
                \texttt{\tiny{[GM$|$GM]}}
                \texttt{Penalizácia: (4,2) nie je prázdny, ale obsahuje '─'.} \\
\\ 
\texttt{Spoločne ste nazbierali 1 z 12 pokút.} \\
\texttt{Skúste to prosím znova!} \\
            }
        }
    }
    & & \\ \\

    \theutterance \stepcounter{utterance}  
    & & \multicolumn{4}{p{0.6\linewidth}}{
        \cellcolor[rgb]{0.9,0.9,0.9}{
            \makecell[{{p{\linewidth}}}]{
                \texttt{\tiny{[GM$\rangle$P2]}}
                \texttt{Penalizácia: (4,2) nie je prázdny, ale obsahuje '─'.} \\
\\ 
\texttt{Spoločne ste nazbierali 1 z 12 pokút.} \\
\texttt{Skúste to prosím znova!} \\
            }
        }
    }
    & \\ \\

    \theutterance \stepcounter{utterance}  
    & & & \multicolumn{4}{p{0.6\linewidth}}{
        \cellcolor[rgb]{0.9,0.9,0.9}{
            \makecell[{{p{\linewidth}}}]{
                \texttt{\tiny{[GM$\langle$P2]}}
                \texttt{POHYB: W, (4, 1)} \\
            }
        }
    }
     \\ \\

    \theutterance \stepcounter{utterance}  
    & & & \multicolumn{2}{p{0.3\linewidth}}{
        \cellcolor[rgb]{0.9,0.9,0.9}{
            \makecell[{{p{\linewidth}}}]{
                \texttt{\tiny{[GM$|$GM]}}
                \texttt{POHYB: W, (4, 1)} \\
            }
        }
    }
    & & \\ \\

    \theutterance \stepcounter{utterance}  
    & & & \multicolumn{2}{p{0.3\linewidth}}{
        \cellcolor[rgb]{0.9,0.9,0.9}{
            \makecell[{{p{\linewidth}}}]{
                \texttt{\tiny{[GM$|$GM]}}
                \texttt{Penalizácia: (4,1) nie je prázdny, ale obsahuje 'W'.} \\
\\ 
\texttt{Spoločne ste nazbierali 2 z 12 pokút.} \\
\texttt{Skúste to prosím znova!} \\
            }
        }
    }
    & & \\ \\

    \theutterance \stepcounter{utterance}  
    & & \multicolumn{4}{p{0.6\linewidth}}{
        \cellcolor[rgb]{0.9,0.9,0.9}{
            \makecell[{{p{\linewidth}}}]{
                \texttt{\tiny{[GM$\rangle$P2]}}
                \texttt{Penalizácia: (4,1) nie je prázdny, ale obsahuje 'W'.} \\
\\ 
\texttt{Spoločne ste nazbierali 2 z 12 pokút.} \\
\texttt{Skúste to prosím znova!} \\
            }
        }
    }
    & \\ \\

    \theutterance \stepcounter{utterance}  
    & & & \multicolumn{4}{p{0.6\linewidth}}{
        \cellcolor[rgb]{0.9,0.9,0.9}{
            \makecell[{{p{\linewidth}}}]{
                \texttt{\tiny{[GM$\langle$P2]}}
                \texttt{POHYB: W, (5, 1)} \\
            }
        }
    }
     \\ \\

    \theutterance \stepcounter{utterance}  
    & & & \multicolumn{2}{p{0.3\linewidth}}{
        \cellcolor[rgb]{0.9,0.9,0.9}{
            \makecell[{{p{\linewidth}}}]{
                \texttt{\tiny{[GM$|$GM]}}
                \texttt{POHYB: W, (5, 1)} \\
            }
        }
    }
    & & \\ \\

    \theutterance \stepcounter{utterance}  
    & & & \multicolumn{2}{p{0.3\linewidth}}{
        \cellcolor[rgb]{0.9,0.9,0.9}{
            \makecell[{{p{\linewidth}}}]{
                \texttt{\tiny{[GM$|$GM]}}
                \texttt{W moved to (5, 1)} \\
            }
        }
    }
    & & \\ \\

    \theutterance \stepcounter{utterance}  
    & & \multicolumn{4}{p{0.6\linewidth}}{
        \cellcolor[rgb]{0.9,0.9,0.9}{
            \makecell[{{p{\linewidth}}}]{
                \texttt{\tiny{[P1$\langle$GM]}}
                \texttt{Úspešne presunutý 'T' na (4,3). Vaša aktualizovaná mriežka vyzerá takto:} \\
\texttt{```} \\
\texttt{ 1234567} \\
\texttt{╔══╤══╤═╗} \\
\texttt{║◌◌│◌◌│◌║ 1} \\
\texttt{╟──┼──┴─╢ 2} \\
\texttt{║WI│T◌◌◌║ 3} \\
\texttt{╟──┘◌◌◌◌║ 4} \\
\texttt{║◌◌◌H◌◌C║ 5} \\
\texttt{║◌◌┌─┬──╢ 6} \\
\texttt{║◌◌│◌│◌◌║ 7} \\
\texttt{╚══╧═╧══╝} \\
\texttt{```} \\
\texttt{Práve hráte kolo 3 z maximálne 20.} \\
\texttt{Spoločne ste nazbierali 2 z 12 pokút.} \\
\texttt{Druhý hráč presunul objekt na svojej mriežke.} \\
\\ 
\\ 
\texttt{Aký je váš ďalší príkaz?} \\
            }
        }
    }
    & \\ \\

    \theutterance \stepcounter{utterance}  
    & \multicolumn{4}{p{0.6\linewidth}}{
        \cellcolor[rgb]{0.9,0.9,0.9}{
            \makecell[{{p{\linewidth}}}]{
                \texttt{\tiny{[P1$\rangle$GM]}}
                \texttt{POHYB: W, (4, 2)} \\
            }
        }
    }
    & & \\ \\

    \theutterance \stepcounter{utterance}  
    & & & \multicolumn{2}{p{0.3\linewidth}}{
        \cellcolor[rgb]{0.9,0.9,0.9}{
            \makecell[{{p{\linewidth}}}]{
                \texttt{\tiny{[GM$|$GM]}}
                \texttt{POHYB: W, (4, 2)} \\
            }
        }
    }
    & & \\ \\

    \theutterance \stepcounter{utterance}  
    & & & \multicolumn{2}{p{0.3\linewidth}}{
        \cellcolor[rgb]{0.9,0.9,0.9}{
            \makecell[{{p{\linewidth}}}]{
                \texttt{\tiny{[GM$|$GM]}}
                \texttt{Penalizácia: (4,2) nie je prázdny, ale obsahuje '─'.} \\
\\ 
\texttt{Spoločne ste nazbierali 3 z 12 pokút.} \\
\texttt{Skúste to prosím znova!} \\
            }
        }
    }
    & & \\ \\

    \theutterance \stepcounter{utterance}  
    & & \multicolumn{4}{p{0.6\linewidth}}{
        \cellcolor[rgb]{0.9,0.9,0.9}{
            \makecell[{{p{\linewidth}}}]{
                \texttt{\tiny{[P1$\langle$GM]}}
                \texttt{Penalizácia: (4,2) nie je prázdny, ale obsahuje '─'.} \\
\\ 
\texttt{Spoločne ste nazbierali 3 z 12 pokút.} \\
\texttt{Skúste to prosím znova!} \\
            }
        }
    }
    & \\ \\

    \theutterance \stepcounter{utterance}  
    & \multicolumn{4}{p{0.6\linewidth}}{
        \cellcolor[rgb]{0.9,0.9,0.9}{
            \makecell[{{p{\linewidth}}}]{
                \texttt{\tiny{[P1$\rangle$GM]}}
                \texttt{POHYB: W, (3, 2)} \\
            }
        }
    }
    & & \\ \\

    \theutterance \stepcounter{utterance}  
    & & & \multicolumn{2}{p{0.3\linewidth}}{
        \cellcolor[rgb]{0.9,0.9,0.9}{
            \makecell[{{p{\linewidth}}}]{
                \texttt{\tiny{[GM$|$GM]}}
                \texttt{POHYB: W, (3, 2)} \\
            }
        }
    }
    & & \\ \\

    \theutterance \stepcounter{utterance}  
    & & & \multicolumn{2}{p{0.3\linewidth}}{
        \cellcolor[rgb]{0.9,0.9,0.9}{
            \makecell[{{p{\linewidth}}}]{
                \texttt{\tiny{[GM$|$GM]}}
                \texttt{Penalizácia: (3,2) nie je prázdny, ale obsahuje '┼'.} \\
\\ 
\texttt{Spoločne ste nazbierali 4 z 12 pokút.} \\
\texttt{Skúste to prosím znova!} \\
            }
        }
    }
    & & \\ \\

    \theutterance \stepcounter{utterance}  
    & & \multicolumn{4}{p{0.6\linewidth}}{
        \cellcolor[rgb]{0.9,0.9,0.9}{
            \makecell[{{p{\linewidth}}}]{
                \texttt{\tiny{[P1$\langle$GM]}}
                \texttt{Penalizácia: (3,2) nie je prázdny, ale obsahuje '┼'.} \\
\\ 
\texttt{Spoločne ste nazbierali 4 z 12 pokút.} \\
\texttt{Skúste to prosím znova!} \\
            }
        }
    }
    & \\ \\

    \theutterance \stepcounter{utterance}  
    & \multicolumn{4}{p{0.6\linewidth}}{
        \cellcolor[rgb]{0.9,0.9,0.9}{
            \makecell[{{p{\linewidth}}}]{
                \texttt{\tiny{[P1$\rangle$GM]}}
                \texttt{POHYB: W, (1, 2)} \\
            }
        }
    }
    & & \\ \\

    \theutterance \stepcounter{utterance}  
    & & & \multicolumn{2}{p{0.3\linewidth}}{
        \cellcolor[rgb]{0.9,0.9,0.9}{
            \makecell[{{p{\linewidth}}}]{
                \texttt{\tiny{[GM$|$GM]}}
                \texttt{POHYB: W, (1, 2)} \\
            }
        }
    }
    & & \\ \\

    \theutterance \stepcounter{utterance}  
    & & & \multicolumn{2}{p{0.3\linewidth}}{
        \cellcolor[rgb]{0.9,0.9,0.9}{
            \makecell[{{p{\linewidth}}}]{
                \texttt{\tiny{[GM$|$GM]}}
                \texttt{Penalizácia: (1,2) nie je prázdny, ale obsahuje '─'.} \\
\\ 
\texttt{Spoločne ste nazbierali 5 z 12 pokút.} \\
\texttt{Skúste to prosím znova!} \\
            }
        }
    }
    & & \\ \\

    \theutterance \stepcounter{utterance}  
    & & \multicolumn{4}{p{0.6\linewidth}}{
        \cellcolor[rgb]{0.9,0.9,0.9}{
            \makecell[{{p{\linewidth}}}]{
                \texttt{\tiny{[P1$\langle$GM]}}
                \texttt{Penalizácia: (1,2) nie je prázdny, ale obsahuje '─'.} \\
\\ 
\texttt{Spoločne ste nazbierali 5 z 12 pokút.} \\
\texttt{Skúste to prosím znova!} \\
            }
        }
    }
    & \\ \\

    \theutterance \stepcounter{utterance}  
    & \multicolumn{4}{p{0.6\linewidth}}{
        \cellcolor[rgb]{0.9,0.9,0.9}{
            \makecell[{{p{\linewidth}}}]{
                \texttt{\tiny{[P1$\rangle$GM]}}
                \texttt{POHYB: W, (1, 4)} \\
            }
        }
    }
    & & \\ \\

    \theutterance \stepcounter{utterance}  
    & & & \multicolumn{2}{p{0.3\linewidth}}{
        \cellcolor[rgb]{0.9,0.9,0.9}{
            \makecell[{{p{\linewidth}}}]{
                \texttt{\tiny{[GM$|$GM]}}
                \texttt{POHYB: W, (1, 4)} \\
            }
        }
    }
    & & \\ \\

    \theutterance \stepcounter{utterance}  
    & & & \multicolumn{2}{p{0.3\linewidth}}{
        \cellcolor[rgb]{0.9,0.9,0.9}{
            \makecell[{{p{\linewidth}}}]{
                \texttt{\tiny{[GM$|$GM]}}
                \texttt{Penalizácia: (1,4) nie je prázdny, ale obsahuje '─'.} \\
\\ 
\texttt{Spoločne ste nazbierali 6 z 12 pokút.} \\
\texttt{Skúste to prosím znova!} \\
            }
        }
    }
    & & \\ \\

    \theutterance \stepcounter{utterance}  
    & & \multicolumn{4}{p{0.6\linewidth}}{
        \cellcolor[rgb]{0.9,0.9,0.9}{
            \makecell[{{p{\linewidth}}}]{
                \texttt{\tiny{[P1$\langle$GM]}}
                \texttt{Penalizácia: (1,4) nie je prázdny, ale obsahuje '─'.} \\
\\ 
\texttt{Spoločne ste nazbierali 6 z 12 pokút.} \\
\texttt{Skúste to prosím znova!} \\
            }
        }
    }
    & \\ \\

    \theutterance \stepcounter{utterance}  
    & \multicolumn{4}{p{0.6\linewidth}}{
        \cellcolor[rgb]{0.9,0.9,0.9}{
            \makecell[{{p{\linewidth}}}]{
                \texttt{\tiny{[P1$\rangle$GM]}}
                \texttt{POHYB: W, (1, 1)} \\
            }
        }
    }
    & & \\ \\

    \theutterance \stepcounter{utterance}  
    & & & \multicolumn{2}{p{0.3\linewidth}}{
        \cellcolor[rgb]{0.9,0.9,0.9}{
            \makecell[{{p{\linewidth}}}]{
                \texttt{\tiny{[GM$|$GM]}}
                \texttt{POHYB: W, (1, 1)} \\
            }
        }
    }
    & & \\ \\

    \theutterance \stepcounter{utterance}  
    & & & \multicolumn{2}{p{0.3\linewidth}}{
        \cellcolor[rgb]{0.9,0.9,0.9}{
            \makecell[{{p{\linewidth}}}]{
                \texttt{\tiny{[GM$|$GM]}}
                \texttt{W moved to (1, 1)} \\
            }
        }
    }
    & & \\ \\

    \theutterance \stepcounter{utterance}  
    & & \multicolumn{4}{p{0.6\linewidth}}{
        \cellcolor[rgb]{0.9,0.9,0.9}{
            \makecell[{{p{\linewidth}}}]{
                \texttt{\tiny{[GM$\rangle$P2]}}
                \texttt{Úspešne presunutý 'W' na (5,1). Vaša aktualizovaná mriežka vyzerá takto:} \\
\texttt{```} \\
\texttt{ 1234567} \\
\texttt{╔══╤══╤═╗} \\
\texttt{║◌◌│◌W│I║ 1} \\
\texttt{╟──┼──┴─╢ 2} \\
\texttt{║◌◌│T◌◌◌║ 3} \\
\texttt{╟──┘◌◌◌◌║ 4} \\
\texttt{║H◌◌◌◌◌◌║ 5} \\
\texttt{║◌◌┌─┬──╢ 6} \\
\texttt{║C◌│◌│◌◌║ 7} \\
\texttt{╚══╧═╧══╝} \\
\texttt{```} \\
\texttt{Práve hráte kolo 4 z maximálne 20.} \\
\texttt{Spoločne ste nazbierali 6 z 12 pokút.} \\
\texttt{Druhý hráč presunul objekt na svojej mriežke.} \\
\\ 
\\ 
\texttt{Aký je váš ďalší príkaz?} \\
            }
        }
    }
    & \\ \\

    \theutterance \stepcounter{utterance}  
    & & & \multicolumn{4}{p{0.6\linewidth}}{
        \cellcolor[rgb]{0.9,0.9,0.9}{
            \makecell[{{p{\linewidth}}}]{
                \texttt{\tiny{[GM$\langle$P2]}}
                \texttt{POHYB: I, (5, 1)} \\
            }
        }
    }
     \\ \\

    \theutterance \stepcounter{utterance}  
    & & & \multicolumn{2}{p{0.3\linewidth}}{
        \cellcolor[rgb]{0.9,0.9,0.9}{
            \makecell[{{p{\linewidth}}}]{
                \texttt{\tiny{[GM$|$GM]}}
                \texttt{POHYB: I, (5, 1)} \\
            }
        }
    }
    & & \\ \\

    \theutterance \stepcounter{utterance}  
    & & & \multicolumn{2}{p{0.3\linewidth}}{
        \cellcolor[rgb]{0.9,0.9,0.9}{
            \makecell[{{p{\linewidth}}}]{
                \texttt{\tiny{[GM$|$GM]}}
                \texttt{Penalizácia: (5,1) nie je prázdny, ale obsahuje 'W'.} \\
\\ 
\texttt{Spoločne ste nazbierali 7 z 12 pokút.} \\
\texttt{Skúste to prosím znova!} \\
            }
        }
    }
    & & \\ \\

    \theutterance \stepcounter{utterance}  
    & & \multicolumn{4}{p{0.6\linewidth}}{
        \cellcolor[rgb]{0.9,0.9,0.9}{
            \makecell[{{p{\linewidth}}}]{
                \texttt{\tiny{[GM$\rangle$P2]}}
                \texttt{Penalizácia: (5,1) nie je prázdny, ale obsahuje 'W'.} \\
\\ 
\texttt{Spoločne ste nazbierali 7 z 12 pokút.} \\
\texttt{Skúste to prosím znova!} \\
            }
        }
    }
    & \\ \\

    \theutterance \stepcounter{utterance}  
    & & & \multicolumn{4}{p{0.6\linewidth}}{
        \cellcolor[rgb]{0.9,0.9,0.9}{
            \makecell[{{p{\linewidth}}}]{
                \texttt{\tiny{[GM$\langle$P2]}}
                \texttt{POHYB: I, (6, 1)} \\
            }
        }
    }
     \\ \\

    \theutterance \stepcounter{utterance}  
    & & & \multicolumn{2}{p{0.3\linewidth}}{
        \cellcolor[rgb]{0.9,0.9,0.9}{
            \makecell[{{p{\linewidth}}}]{
                \texttt{\tiny{[GM$|$GM]}}
                \texttt{POHYB: I, (6, 1)} \\
            }
        }
    }
    & & \\ \\

    \theutterance \stepcounter{utterance}  
    & & & \multicolumn{2}{p{0.3\linewidth}}{
        \cellcolor[rgb]{0.9,0.9,0.9}{
            \makecell[{{p{\linewidth}}}]{
                \texttt{\tiny{[GM$|$GM]}}
                \texttt{Penalizácia: (6,1) nie je prázdny, ale obsahuje '│'.} \\
\\ 
\texttt{Spoločne ste nazbierali 8 z 12 pokút.} \\
\texttt{Skúste to prosím znova!} \\
            }
        }
    }
    & & \\ \\

    \theutterance \stepcounter{utterance}  
    & & \multicolumn{4}{p{0.6\linewidth}}{
        \cellcolor[rgb]{0.9,0.9,0.9}{
            \makecell[{{p{\linewidth}}}]{
                \texttt{\tiny{[GM$\rangle$P2]}}
                \texttt{Penalizácia: (6,1) nie je prázdny, ale obsahuje '│'.} \\
\\ 
\texttt{Spoločne ste nazbierali 8 z 12 pokút.} \\
\texttt{Skúste to prosím znova!} \\
            }
        }
    }
    & \\ \\

    \theutterance \stepcounter{utterance}  
    & & & \multicolumn{4}{p{0.6\linewidth}}{
        \cellcolor[rgb]{0.9,0.9,0.9}{
            \makecell[{{p{\linewidth}}}]{
                \texttt{\tiny{[GM$\langle$P2]}}
                \texttt{POHYB: I, (7, 1)} \\
            }
        }
    }
     \\ \\

    \theutterance \stepcounter{utterance}  
    & & & \multicolumn{2}{p{0.3\linewidth}}{
        \cellcolor[rgb]{0.9,0.9,0.9}{
            \makecell[{{p{\linewidth}}}]{
                \texttt{\tiny{[GM$|$GM]}}
                \texttt{POHYB: I, (7, 1)} \\
            }
        }
    }
    & & \\ \\

    \theutterance \stepcounter{utterance}  
    & & & \multicolumn{2}{p{0.3\linewidth}}{
        \cellcolor[rgb]{0.9,0.9,0.9}{
            \makecell[{{p{\linewidth}}}]{
                \texttt{\tiny{[GM$|$GM]}}
                \texttt{Penalizácia: (7,1) nie je prázdny, ale obsahuje 'I'.} \\
\\ 
\texttt{Spoločne ste nazbierali 9 z 12 pokút.} \\
\texttt{Skúste to prosím znova!} \\
            }
        }
    }
    & & \\ \\

    \theutterance \stepcounter{utterance}  
    & & \multicolumn{4}{p{0.6\linewidth}}{
        \cellcolor[rgb]{0.9,0.9,0.9}{
            \makecell[{{p{\linewidth}}}]{
                \texttt{\tiny{[GM$\rangle$P2]}}
                \texttt{Penalizácia: (7,1) nie je prázdny, ale obsahuje 'I'.} \\
\\ 
\texttt{Spoločne ste nazbierali 9 z 12 pokút.} \\
\texttt{Skúste to prosím znova!} \\
            }
        }
    }
    & \\ \\

    \theutterance \stepcounter{utterance}  
    & & & \multicolumn{4}{p{0.6\linewidth}}{
        \cellcolor[rgb]{0.9,0.9,0.9}{
            \makecell[{{p{\linewidth}}}]{
                \texttt{\tiny{[GM$\langle$P2]}}
                \texttt{POHYB: I, (6, 1)} \\
\\ 
\texttt{*(Poznámka: Musím si uvedomiť, že v mriežke sú stĺpce 1, 2, 3, 4, 5, 6, 7. Podľa vizuálneho znázornenia mriežky sú bunky v riadku 1 na pozíciách (1,1), (2,1), (3,1), (4,1), (5,1), (6,1), (7,1). Avšak mriežka obsahuje aj čiary. Pozorne sa pozriem na mriežku: (7,1) je 'I', (6,1) je '│', (5,1) je 'W', (4,1) je '◌', (3,1) je '│', (2,1) je '◌', (1,1) je '◌'. Moja predchádzajúca chyba bola snaha presunúť I na (6,1) alebo (7,1) čo sú obsadené. Skúsím nájsť prázdnu bunku v riadku 1. Podľa mriežky sú prázdne bunky (1,1), (2,1) a (4,1). Skúsime presunúť I na (4,1).)*} \\
\\ 
\texttt{POHYB: I, (4, 1)} \\
            }
        }
    }
     \\ \\

    \theutterance \stepcounter{utterance}  
    & & & \multicolumn{2}{p{0.3\linewidth}}{
        \cellcolor[rgb]{0.9,0.9,0.9}{
            \makecell[{{p{\linewidth}}}]{
                \texttt{\tiny{[GM$|$GM]}}
                \texttt{POHYB: I, (6, 1)} \\
\\ 
\texttt{*(Poznámka: Musím si uvedomiť, že v mriežke sú stĺpce 1, 2, 3, 4, 5, 6, 7. Podľa vizuálneho znázornenia mriežky sú bunky v riadku 1 na pozíciách (1,1), (2,1), (3,1), (4,1), (5,1), (6,1), (7,1). Avšak mriežka obsahuje aj čiary. Pozorne sa pozriem na mriežku: (7,1) je 'I', (6,1) je '│', (5,1) je 'W', (4,1) je '◌', (3,1) je '│', (2,1) je '◌', (1,1) je '◌'. Moja predchádzajúca chyba bola snaha presunúť I na (6,1) alebo (7,1) čo sú obsadené. Skúsím nájsť prázdnu bunku v riadku 1. Podľa mriežky sú prázdne bunky (1,1), (2,1) a (4,1). Skúsime presunúť I na (4,1).)*} \\
\\ 
\texttt{POHYB: I, (4, 1)} \\
            }
        }
    }
    & & \\ \\

    \theutterance \stepcounter{utterance}  
    & & & \multicolumn{2}{p{0.3\linewidth}}{
        \cellcolor[rgb]{0.9,0.9,0.9}{
            \makecell[{{p{\linewidth}}}]{
                \texttt{\tiny{[GM$|$GM]}}
                \texttt{Invalid format: head is not empty: 'POHYB: I, (6, 1)} \\
\\ 
\texttt{*(Poznámka: Musím si uvedomiť, že v mriežke sú stĺpce 1, 2, 3, 4, 5, 6, 7. Podľa vizuálneho znázornenia mriežky sú bunky v riadku 1 na pozíciách (1,1), (2,1), (3,1), (4,1), (5,1), (6,1), (7,1). Avšak mriežka obsahuje aj čiary. Pozorne sa pozriem na mriežku: (7,1) je 'I', (6,1) je '│', (5,1) je 'W', (4,1) je '◌', (3,1) je '│', (2,1) je '◌', (1,1) je '◌'. Moja predchádzajúca chyba bola snaha presunúť I na (6,1) alebo (7,1) čo sú obsadené. Skúsím nájsť prázdnu bunku v riadku 1. Podľa mriežky sú prázdne bunky (1,1), (2,1) a (4,1). Skúsime presunúť I na (4,1).)*} \\
\\ 
\texttt{'} \\
            }
        }
    }
    & & \\ \\

    \theutterance \stepcounter{utterance}  
    & & \multicolumn{4}{p{0.6\linewidth}}{
        \cellcolor[rgb]{0.9,0.9,0.9}{
            \makecell[{{p{\linewidth}}}]{
                \texttt{\tiny{[GM$\rangle$P2]}}
                \texttt{Trest: Vaša správa nesmie obsahovať nič pred príkazom!} \\
\\ 
\texttt{Uistite sa, že vaša odpoveď obsahuje iba `POVEDZ: <SPRÁVA>` alebo `POHYB: <OBJEKT>, (<X>, <Y>)` a nič iné!} \\
\\ 
\texttt{Spoločne ste nazbierali 10 z 12 pokút.} \\
\texttt{Skúste to prosím znova!} \\
            }
        }
    }
    & \\ \\

    \theutterance \stepcounter{utterance}  
    & & & \multicolumn{4}{p{0.6\linewidth}}{
        \cellcolor[rgb]{0.9,0.9,0.9}{
            \makecell[{{p{\linewidth}}}]{
                \texttt{\tiny{[GM$\langle$P2]}}
                \texttt{POHYB: I, (4, 1)} \\
            }
        }
    }
     \\ \\

    \theutterance \stepcounter{utterance}  
    & & & \multicolumn{2}{p{0.3\linewidth}}{
        \cellcolor[rgb]{0.9,0.9,0.9}{
            \makecell[{{p{\linewidth}}}]{
                \texttt{\tiny{[GM$|$GM]}}
                \texttt{POHYB: I, (4, 1)} \\
            }
        }
    }
    & & \\ \\

    \theutterance \stepcounter{utterance}  
    & & & \multicolumn{2}{p{0.3\linewidth}}{
        \cellcolor[rgb]{0.9,0.9,0.9}{
            \makecell[{{p{\linewidth}}}]{
                \texttt{\tiny{[GM$|$GM]}}
                \texttt{I moved to (4, 1)} \\
            }
        }
    }
    & & \\ \\

    \theutterance \stepcounter{utterance}  
    & & \multicolumn{4}{p{0.6\linewidth}}{
        \cellcolor[rgb]{0.9,0.9,0.9}{
            \makecell[{{p{\linewidth}}}]{
                \texttt{\tiny{[P1$\langle$GM]}}
                \texttt{Úspešne presunutý 'W' na (1,1). Vaša aktualizovaná mriežka vyzerá takto:} \\
\texttt{```} \\
\texttt{ 1234567} \\
\texttt{╔══╤══╤═╗} \\
\texttt{║W◌│◌◌│◌║ 1} \\
\texttt{╟──┼──┴─╢ 2} \\
\texttt{║◌I│T◌◌◌║ 3} \\
\texttt{╟──┘◌◌◌◌║ 4} \\
\texttt{║◌◌◌H◌◌C║ 5} \\
\texttt{║◌◌┌─┬──╢ 6} \\
\texttt{║◌◌│◌│◌◌║ 7} \\
\texttt{╚══╧═╧══╝} \\
\texttt{```} \\
\texttt{Práve hráte kolo 4 z maximálne 20.} \\
\texttt{Spoločne ste nazbierali 10 z 12 pokút.} \\
\texttt{Druhý hráč presunul objekt na svojej mriežke.} \\
\\ 
\\ 
\texttt{Aký je váš ďalší príkaz?} \\
            }
        }
    }
    & \\ \\

    \theutterance \stepcounter{utterance}  
    & \multicolumn{4}{p{0.6\linewidth}}{
        \cellcolor[rgb]{0.9,0.9,0.9}{
            \makecell[{{p{\linewidth}}}]{
                \texttt{\tiny{[P1$\rangle$GM]}}
                \texttt{POHYB: I, (5, 2)} \\
            }
        }
    }
    & & \\ \\

    \theutterance \stepcounter{utterance}  
    & & & \multicolumn{2}{p{0.3\linewidth}}{
        \cellcolor[rgb]{0.9,0.9,0.9}{
            \makecell[{{p{\linewidth}}}]{
                \texttt{\tiny{[GM$|$GM]}}
                \texttt{POHYB: I, (5, 2)} \\
            }
        }
    }
    & & \\ \\

    \theutterance \stepcounter{utterance}  
    & & & \multicolumn{2}{p{0.3\linewidth}}{
        \cellcolor[rgb]{0.9,0.9,0.9}{
            \makecell[{{p{\linewidth}}}]{
                \texttt{\tiny{[GM$|$GM]}}
                \texttt{Penalizácia: (5,2) nie je prázdny, ale obsahuje '─'.} \\
\\ 
\texttt{Spoločne ste nazbierali 11 z 12 pokút.} \\
\texttt{Skúste to prosím znova!} \\
            }
        }
    }
    & & \\ \\

    \theutterance \stepcounter{utterance}  
    & & \multicolumn{4}{p{0.6\linewidth}}{
        \cellcolor[rgb]{0.9,0.9,0.9}{
            \makecell[{{p{\linewidth}}}]{
                \texttt{\tiny{[P1$\langle$GM]}}
                \texttt{Penalizácia: (5,2) nie je prázdny, ale obsahuje '─'.} \\
\\ 
\texttt{Spoločne ste nazbierali 11 z 12 pokút.} \\
\texttt{Skúste to prosím znova!} \\
            }
        }
    }
    & \\ \\

    \theutterance \stepcounter{utterance}  
    & \multicolumn{4}{p{0.6\linewidth}}{
        \cellcolor[rgb]{0.9,0.9,0.9}{
            \makecell[{{p{\linewidth}}}]{
                \texttt{\tiny{[P1$\rangle$GM]}}
                \texttt{POHYB: I, (5, 3)} \\
            }
        }
    }
    & & \\ \\

    \theutterance \stepcounter{utterance}  
    & & & \multicolumn{2}{p{0.3\linewidth}}{
        \cellcolor[rgb]{0.9,0.9,0.9}{
            \makecell[{{p{\linewidth}}}]{
                \texttt{\tiny{[GM$|$GM]}}
                \texttt{POHYB: I, (5, 3)} \\
            }
        }
    }
    & & \\ \\

    \theutterance \stepcounter{utterance}  
    & & & \multicolumn{2}{p{0.3\linewidth}}{
        \cellcolor[rgb]{0.9,0.9,0.9}{
            \makecell[{{p{\linewidth}}}]{
                \texttt{\tiny{[GM$|$GM]}}
                \texttt{I moved to (5, 3)} \\
            }
        }
    }
    & & \\ \\

    \theutterance \stepcounter{utterance}  
    & & \multicolumn{4}{p{0.6\linewidth}}{
        \cellcolor[rgb]{0.9,0.9,0.9}{
            \makecell[{{p{\linewidth}}}]{
                \texttt{\tiny{[GM$\rangle$P2]}}
                \texttt{Úspešne presunutý 'I' na (4,1). Vaša aktualizovaná mriežka vyzerá takto:} \\
\texttt{```} \\
\texttt{ 1234567} \\
\texttt{╔══╤══╤═╗} \\
\texttt{║◌◌│IW│◌║ 1} \\
\texttt{╟──┼──┴─╢ 2} \\
\texttt{║◌◌│T◌◌◌║ 3} \\
\texttt{╟──┘◌◌◌◌║ 4} \\
\texttt{║H◌◌◌◌◌◌║ 5} \\
\texttt{║◌◌┌─┬──╢ 6} \\
\texttt{║C◌│◌│◌◌║ 7} \\
\texttt{╚══╧═╧══╝} \\
\texttt{```} \\
\texttt{Práve hráte kolo 5 z maximálne 20.} \\
\texttt{Spoločne ste nazbierali 11 z 12 pokút.} \\
\texttt{Druhý hráč presunul objekt na svojej mriežke.} \\
\\ 
\\ 
\texttt{Aký je váš ďalší príkaz?} \\
            }
        }
    }
    & \\ \\

    \theutterance \stepcounter{utterance}  
    & & & \multicolumn{4}{p{0.6\linewidth}}{
        \cellcolor[rgb]{0.9,0.9,0.9}{
            \makecell[{{p{\linewidth}}}]{
                \texttt{\tiny{[GM$\langle$P2]}}
                \texttt{POHYB: I, (5, 2)} \\
            }
        }
    }
     \\ \\

    \theutterance \stepcounter{utterance}  
    & & & \multicolumn{2}{p{0.3\linewidth}}{
        \cellcolor[rgb]{0.9,0.9,0.9}{
            \makecell[{{p{\linewidth}}}]{
                \texttt{\tiny{[GM$|$GM]}}
                \texttt{POHYB: I, (5, 2)} \\
            }
        }
    }
    & & \\ \\

    \theutterance \stepcounter{utterance}  
    & & & \multicolumn{2}{p{0.3\linewidth}}{
        \cellcolor[rgb]{0.9,0.9,0.9}{
            \makecell[{{p{\linewidth}}}]{
                \texttt{\tiny{[GM$|$GM]}}
                \texttt{Penalizácia: (5,2) nie je prázdny, ale obsahuje '─'.} \\
\\ 
\texttt{Spoločne ste nazbierali 12 z 12 pokút.} \\
\texttt{Skúste to prosím znova!} \\
            }
        }
    }
    & & \\ \\

    \theutterance \stepcounter{utterance}  
    & & \multicolumn{4}{p{0.6\linewidth}}{
        \cellcolor[rgb]{0.9,0.9,0.9}{
            \makecell[{{p{\linewidth}}}]{
                \texttt{\tiny{[GM$\rangle$P2]}}
                \texttt{Penalizácia: (5,2) nie je prázdny, ale obsahuje '─'.} \\
\\ 
\texttt{Spoločne ste nazbierali 12 z 12 pokút.} \\
\texttt{Skúste to prosím znova!} \\
            }
        }
    }
    & \\ \\

    \theutterance \stepcounter{utterance}  
    & & & \multicolumn{4}{p{0.6\linewidth}}{
        \cellcolor[rgb]{0.9,0.9,0.9}{
            \makecell[{{p{\linewidth}}}]{
                \texttt{\tiny{[GM$\langle$P2]}}
                \texttt{POHYB: I, (6, 1)} \\
\\ 
\texttt{*(Opätovne kontrolujem mriežku: Riadok 1 je `║◌◌│IW│◌║`. Pozície sú: (1,1)=◌, (2,1)=◌, (3,1)=│, (4,1)=I, (5,1)=W, (6,1)=│, (7,1)=◌. Moja predchádzajúca chyba bola, že som sa pokúsil presunúť na (5,2), čo je čiaru. Teraz sa pokúsim presunúť I na (7,1), čo je prázdna bunka v riadku 1.)*} \\
\\ 
\texttt{POHYB: I, (7, 1)} \\
            }
        }
    }
     \\ \\

    \theutterance \stepcounter{utterance}  
    & & & \multicolumn{2}{p{0.3\linewidth}}{
        \cellcolor[rgb]{0.9,0.9,0.9}{
            \makecell[{{p{\linewidth}}}]{
                \texttt{\tiny{[GM$|$GM]}}
                \texttt{POHYB: I, (6, 1)} \\
\\ 
\texttt{*(Opätovne kontrolujem mriežku: Riadok 1 je `║◌◌│IW│◌║`. Pozície sú: (1,1)=◌, (2,1)=◌, (3,1)=│, (4,1)=I, (5,1)=W, (6,1)=│, (7,1)=◌. Moja predchádzajúca chyba bola, že som sa pokúsil presunúť na (5,2), čo je čiaru. Teraz sa pokúsim presunúť I na (7,1), čo je prázdna bunka v riadku 1.)*} \\
\\ 
\texttt{POHYB: I, (7, 1)} \\
            }
        }
    }
    & & \\ \\

    \theutterance \stepcounter{utterance}  
    & & & \multicolumn{2}{p{0.3\linewidth}}{
        \cellcolor[rgb]{0.9,0.9,0.9}{
            \makecell[{{p{\linewidth}}}]{
                \texttt{\tiny{[GM$|$GM]}}
                \texttt{Invalid format: head is not empty: 'POHYB: I, (6, 1)} \\
\\ 
\texttt{*(Opätovne kontrolujem mriežku: Riadok 1 je ║◌◌│IW│◌║. Pozície sú: (1,1)=◌, (2,1)=◌, (3,1)=│, (4,1)=I, (5,1)=W, (6,1)=│, (7,1)=◌. Moja predchádzajúca chyba bola, že som sa pokúsil presunúť na (5,2), čo je čiaru. Teraz sa pokúsim presunúť I na (7,1), čo je prázdna bunka v riadku 1.)*} \\
\\ 
\texttt{'} \\
            }
        }
    }
    & & \\ \\

    \theutterance \stepcounter{utterance}  
    & & & \multicolumn{2}{p{0.3\linewidth}}{
        \cellcolor[rgb]{0.9,0.9,0.9}{
            \makecell[{{p{\linewidth}}}]{
                \texttt{\tiny{[GM$|$GM]}}
                \texttt{Maximum number of penalties exceeded} \\
            }
        }
    }
    & & \\ \\

    \theutterance \stepcounter{utterance}  
    & & & \multicolumn{2}{p{0.3\linewidth}}{
        \cellcolor[rgb]{0.9,0.9,0.9}{
            \makecell[{{p{\linewidth}}}]{
                \texttt{\tiny{[GM$|$GM]}}
                \texttt{Initial states:} \\
\texttt{```} \\
\texttt{Player 1:} \\
\texttt{ 1234567} \\
\texttt{╔══╤══╤═╗} \\
\texttt{║◌◌│T◌│◌║ 1} \\
\texttt{╟──┼──┴─╢ 2} \\
\texttt{║WI│◌◌◌◌║ 3} \\
\texttt{╟──┘◌◌◌◌║ 4} \\
\texttt{║◌◌◌H◌◌C║ 5} \\
\texttt{║◌◌┌─┬──╢ 6} \\
\texttt{║◌◌│◌│◌◌║ 7} \\
\texttt{╚══╧═╧══╝} \\
\\ 
\\ 
\texttt{Player 2:} \\
\texttt{ 1234567} \\
\texttt{╔══╤══╤═╗} \\
\texttt{║◌◌│W◌│I║ 1} \\
\texttt{╟──┼──┴─╢ 2} \\
\texttt{║◌◌│◌◌◌◌║ 3} \\
\texttt{╟──┘T◌◌◌║ 4} \\
\texttt{║H◌◌◌◌◌◌║ 5} \\
\texttt{║◌◌┌─┬──╢ 6} \\
\texttt{║C◌│◌│◌◌║ 7} \\
\texttt{╚══╧═╧══╝} \\
\\ 
\texttt{```} \\
            }
        }
    }
    & & \\ \\

    \theutterance \stepcounter{utterance}  
    & & & \multicolumn{2}{p{0.3\linewidth}}{
        \cellcolor[rgb]{0.9,0.9,0.9}{
            \makecell[{{p{\linewidth}}}]{
                \texttt{\tiny{[GM$|$GM]}}
                \texttt{End states:} \\
\texttt{```} \\
\texttt{Player 1:} \\
\texttt{ 1234567} \\
\texttt{╔══╤══╤═╗} \\
\texttt{║W◌│◌◌│◌║ 1} \\
\texttt{╟──┼──┴─╢ 2} \\
\texttt{║◌◌│TI◌◌║ 3} \\
\texttt{╟──┘◌◌◌◌║ 4} \\
\texttt{║◌◌◌H◌◌C║ 5} \\
\texttt{║◌◌┌─┬──╢ 6} \\
\texttt{║◌◌│◌│◌◌║ 7} \\
\texttt{╚══╧═╧══╝} \\
\\ 
\\ 
\texttt{Player 2:} \\
\texttt{ 1234567} \\
\texttt{╔══╤══╤═╗} \\
\texttt{║◌◌│IW│◌║ 1} \\
\texttt{╟──┼──┴─╢ 2} \\
\texttt{║◌◌│T◌◌◌║ 3} \\
\texttt{╟──┘◌◌◌◌║ 4} \\
\texttt{║H◌◌◌◌◌◌║ 5} \\
\texttt{║◌◌┌─┬──╢ 6} \\
\texttt{║C◌│◌│◌◌║ 7} \\
\texttt{╚══╧═╧══╝} \\
\\ 
\texttt{```} \\
            }
        }
    }
    & & \\ \\

    \theutterance \stepcounter{utterance}  
    & & & \multicolumn{2}{p{0.3\linewidth}}{
        \cellcolor[rgb]{0.9,0.9,0.9}{
            \makecell[{{p{\linewidth}}}]{
                \texttt{\tiny{[GM$|$GM]}}
                \texttt{* success: False} \\
\texttt{* lose: False} \\
\texttt{* aborted: True} \\
\texttt{{-}{-}{-}{-}{-}{-}{-}} \\
\texttt{* Shifts: 2.00} \\
\texttt{* Max Shifts: 8.00} \\
\texttt{* Min Shifts: 4.00} \\
\texttt{* End Distance Sum: 15.56} \\
\texttt{* Init Distance Sum: 21.32} \\
\texttt{* Expected Distance Sum: 20.95} \\
\texttt{* Penalties: 13.00} \\
\texttt{* Max Penalties: 12.00} \\
\texttt{* Rounds: 4.00} \\
\texttt{* Max Rounds: 20.00} \\
\texttt{* Object Count: 5.00} \\
            }
        }
    }
    & & \\ \\

\end{supertabular}
}

\end{document}
